% 编译方式：latexmk + XeLaTeX（.latexmkrc 负责把 texmf/ 加入宏包搜索路径）
% !LW recipe = latexmk (xelatex)
% 使用 TeX Live 自带中文字体，保持与模板相同的字体方案
\PassOptionsToClass{fontset=fandol}{ctexbook}
\documentclass[forprint,AutoFakeBold,AutoFakeSlant]{WHUExperiment}
\usepackage{etoolbox}
\usepackage{needspace}
\usepackage{remreset}
\makeatletter
\patchcmd{\maketitle}{武汉大学计算机学院}{武汉大学经济与管理学院}{}{\errmessage{Cover school patch failed}}
\patchcmd{\maketitle}{本科生课程实验报告}{研究生课程研究报告}{}{\errmessage{Cover title patch failed}}
\patchcmd{\maketitle}{\raisebox{-3ex}[0pt]{\makebox[3.5cm][s]{指\hfill 导\hfill 教\hfill 师\hfill:\hfill}} & { {}\raisebox{-3ex}[0pt]{\the\Csupervisor}\hfill{}} & \\[3ex]}{}{}{\errmessage{Supervisor row patch failed}}
% 保留原模板字号，修正加长后的学院名称居中及封面留白
\patchcmd{\maketitle}{\makebox[9cm][s]{武汉大学经济与管理学院}}{\makebox[\textwidth][c]{武汉大学经济与管理学院}}{}{\errmessage{Cover alignment patch failed}}
\patchcmd{\maketitle}{\vspace{2.0cm}}{\vspace{0.6cm}}{}{\errmessage{Cover spacing patch failed}}
\patchcmd{\maketitle}{\vfill\vfill\vfill}{\vspace{1.2cm}}{}{\errmessage{Cover gap patch failed}}
\patchcmd{\maketitle}{cp{3cm}c}{cp{6cm}c}{}{\errmessage{Cover information width patch failed}}
\makeatother
\hypersetup{hypertexnames=false,pdftitle={柏顿·墨基尔选股法的A股复现与延展},pdfauthor={伊克兰·穆提拉},pdfkeywords={墨基尔选股法；持续盈余成长；低市盈率；排名母体；波动率目标；策略复现}}

\begin{document}
\Ccoursename{金融数据分析}
\title{柏顿·墨基尔选股法的A股复现与延展\\[0.5em]{\zihao{3}四条优化路径的收益、风险与机制比较}}
\author{伊克兰·穆提拉}
\CstudentNum{2026281050061}
\Cmajor{金融}
\date{二〇二六年十一月七日}
\pdfbookmark[0]{封面}{title}
\maketitle

\mainmatter
\pagestyle{plain}
\baselineskip=22pt
\raggedbottom
\widowpenalty=10000
\clubpenalty=10000
\displaywidowpenalty=10000
\setlength{\abovedisplayskip}{8pt plus 2pt minus 2pt}
\setlength{\belowdisplayskip}{8pt plus 2pt minus 2pt}
\setlength{\abovedisplayshortskip}{4pt plus 2pt}
\setlength{\belowdisplayshortskip}{6pt plus 2pt minus 2pt}

% 全文表格、图按出现顺序连续编号，不随章重置
\makeatletter
\@removefromreset{figure}{chapter}
\@removefromreset{table}{chapter}
\makeatother
\renewcommand{\thefigure}{\arabic{figure}}
\renewcommand{\thetable}{\arabic{table}}

% !TeX root = ../main.tex

\begin{cnabstract}
本文以申万宏源研究中的柏顿·墨基尔选股法为复现对象，保留基础复现与同规则延展，并建立四条有明确对照关系的优化路径：正PE与初始五年修正、全截面独立排名，以及分别叠加原有仓位管理的两条路径。原始时期的补充检验中，151个月敏感性口径下P1、P2年化分别为38.77\%、40.97\%，高于同窗口基础版36.44\%；两条风控路径以降低收益换取回撤收窄。基础复现实际覆盖2003年4月至2015年6月147个月；完整延展及四路径统一覆盖至2026年9月，共282个月。结果显示，基础延展年化收益为26.17\%，最大回撤为$-$70.11\%；路径一、二年化收益分别为29.64\%和27.38\%，最大回撤分别为$-$68.87\%和$-$64.15\%。路径三、四将最大回撤收窄至$-$32.74\%和$-$34.86\%，但年化收益降至23.19\%和20.85\%。路径一在全期收益上居前，路径三在所设零利率夏普与回撤指标上较优，但分段排序不稳定；2025—2026年追加区间中，基础延展优于四条路径。机制核查表明，排名母体会改变补位池及最终核心保留，风控路径未出现零仓位月份，不能将风险改善归因于已经生效的止损清仓。四路径仍沿用报告期财务映射、月度再等权和费用近似，独立CSMAR重建逐股复现基准，282个月最大收益差约$5\times10^{-16}$；与原报告的差异可定位到比较区间、调仓安排及部分历史缺口。本文据此将其解释为具有可复核计算依据的历史规则比较，而非已证明的可交易超额收益。
\end{cnabstract}
\cnkeywords{墨基尔选股法；持续盈余成长；低市盈率；排名母体；波动率目标；策略复现}

% !TeX root = ../main.tex

\chapter{引言}
\section{研究背景与问题}
以合理价格购买持续成长企业，是价值与成长相结合的选股思路。柏顿·墨基尔成功选股法在A股的量化实现，以历史盈余成长、当前成长排名和相对低市盈率共同限制候选股票。申万宏源研究[1]提供了可复现的规则与回溯收益，适合检验投资理念在本地数据与具体程序中的表现。

复现的难点不只在于得到一条上涨净值曲线。历史记录是否齐全、百分比排名以何种股票为分母、负PE是否参与排序、候选稀缺时怎样构造组合，都会改变最终名单。在此基础上引入多因子得分和动态仓位，还需要区分选股改变与风险暴露调整，否则容易把多项同时发生的变化归因于某一个“有效因子”。

本文首先回答三个问题：第一，基础规则在原复现区间与延展区间的收益和风险是否一致；第二，候选池内排名与全截面独立排名如何改变组合；第三，在两种选股路径上叠加同一仓位规则，能否以可接受的收益代价降低风险。研究同时保留不利结果，不以更高回溯收益作为判断规则是否正确的标准。

\section{文献联系与本文定位}
本文直接复现对象为券商研究[1]。Fama和French[2]的横截面收益研究提示，组合的高收益还需放在市场和风格暴露背景下理解；价值与动量结合的研究[3]为多维选股提供相关背景，但不能直接证明本项目质量、动量与规模等权得分的有效性。本文没有将其他市场的经验结果移植为A股结论，也未额外引入新的定价模型。新增差异核验以原文、数据覆盖及独立计算作为证据，不根据收益率接近程度决定规则。

本文的工作在于建立可核查的规则链条。基础复现与延展保持同一筛选逻辑；路径一在已修正正PE与初始期规则的代码上保留19/20和三因子选择；路径二仅改变两项排名及各自分母；路径三、四分别只改变前两条路径的仓位。成对比较可以明确实现差异的经济幅度，但不等于随机实验或单因子的因果识别。

\section{原研报与基础复现的基本情况对照}
本研究的直接复现对象是申万宏源2015年发布的柏顿·墨基尔成功选股法[1]，比较对象首先应是基础复现B。延展E改变了结束时间，P1—P4还改变了部分选股或仓位规则，均不能直接替代B来判断原策略是否复现。表1列出研究设计，表2列出原文公布值与本项目实际计算值；原文未披露的细节按“未详”保留，不作一致性假定。

\begin{table}[htbp]
\centering
\caption{申万原研报与基础复现的研究设计对照}
\renewcommand{\arraystretch}{1.18}
\small
\setlength{\tabcolsep}{4.5pt}
\begin{tabularx}{\textwidth}{@{}p{2.45cm}LL@{}}
\toprule
比较项目 & 申万原研报披露 & 本研究基础复现B \\
\midrule
研究对象 & 历史全部A股，含后来退市股票 & 历史A股范围；按现有CSMAR估值、财务和收益覆盖形成可计算集合 \\
成长率 & 净利润同比增长率 & 累计总净利润同比；上年同期利润非正或缺失时，不形成有效成长率 \\
估值指标 & 股价/每股收益（TTM） & 月末总市值/总净利润（TTM）；归属与股本口径尚未逐项确认等价 \\
成熟期筛选 & 过去五年持续正成长，成长率前30\%；正PE且低估值30\% & 最近20季齐全且全为正；当期成长率前30\%与估值最低30\%取交集 \\
初始五年 & 历史每期成长率为正，成长率最高前15只 & 先取历史合格与估值条件交集，再按成长率选最多15只 \\
排名与缺失值 & 市场相对排序；负PE是否占名次、缺失历史处理等细节未完全披露 & 成长和PE各自全截面排名；负PE先占位，再要求入选PE为正；缺失不视作正成长 \\
调仓时点 & 每年四次；全市场最晚财报发布日与报告期后90个交易日孰早 & 固定3/6/9/12月末映射当季财报，次月起持有三个月 \\
权重与费用 & 资金平均配置；交易费率0.13\%，逐笔实现未详 & 每月对有收益记录的持仓等权；有持仓季度首月固定扣0.13\% \\
收益与基准 & 日频回测；上证综指 & 个股含红利月收益；上证综指日收益复利为月收益 \\
\bottomrule
\end{tabularx}
\par\smallskip
\begin{minipage}{\textwidth}\small 资料来源：申万原研报[1]第5—6页及本项目基础程序。初始五年的先后筛选顺序、亏损PE的排名占位和历史资料覆盖，原文不能唯一确定；本表如实列出本项目实现。报告期末财务并非当日已公开信息，本研究调仓近似不等价于原文发布日规则。\end{minipage}
\end{table}

\begin{table}[htbp]
\centering
\caption{申万原研报与基础复现的主要结果对照}
\renewcommand{\arraystretch}{1.18}
\small
\setlength{\tabcolsep}{4.5pt}
\begin{tabularx}{\textwidth}{@{}p{2.05cm}LLL@{}}
\toprule
比较项目 & 申万原研报披露 & 本研究基础复现B & 比较说明 \\
\midrule
实际统计区间 & 原文称2003年至2015-07-31；2938个交易日 & 2003-04至2015-06；147个月 & 起点细节及终点未完全对齐 \\
年化收益 & 36.74\% & 38.86\% & 相差2.12个百分点；不是同日同频差额 \\
累计收益 & 4176.87\% & 5481.49\% & 须同时阅读统计区间 \\
最大回撤幅度 & 70.19\% & 70.11\% & 原日频、本月末；相近不证明路径相同 \\
夏普比率 & 0.93 & 1.100（零无风险利率） & 频率及无风险口径未统一 \\
平均持股数 & 7只 & 8.02只（49次信号等权平均） & 原文均值汇总细节未披露 \\
最少/最多持股 & 0/15只 & 0/15只 & 仅比较观测极值，非容量规则相同 \\
空仓情况 & 2段空仓，具体起止见下文 & 3个空仓季度信号，共9个月 & 日历时段和信号数不属同一单位 \\
末期组合 & 2015-05-25调仓；2015-07-31持有8只 & 2015-03-31信号；持有9只至6月 & 共有4只；选股日和观察终点不同 \\
\bottomrule
\end{tabularx}
\par\smallskip
\begin{minipage}{\textwidth}\small 资料来源：原研报[1]第6、8、12页；本研究六方案月收益、逐季持股表及统一指标。最大回撤为便于对照取损失幅度正值，后文绩效表以负值表示回撤。B的CAGR使用147/12年；原文的精确建仓日、年化日数约定及夏普无风险利率未充分披露，原图首个刻度不作建仓日依据。\end{minipage}
\end{table}

收益方面，B的年化38.86\%与原文36.74\%在数量级上接近，但累计收益、调仓日期和末期名单均未完全一致，不能据此宣布精确复现。风险方面，约70\%的峰谷损失是两者共同呈现的显著特征；日频与月末回撤相近只表明本次数值接近，不能证明月内路径或风险控制能力相同。原文夏普0.93与本研究零利率月频夏普1.100也不作严格优劣判断。

持仓方面，原文两段空仓为2009年11月2日至2010年6月29日、2010年9月1日至10月29日；B的空仓信号为2009年9月30日、2010年3月31日、2013年3月29日，对应2009年第四季度、2010年第二季度及2013年第二季度各三个月。两者有部分时间重叠，但空仓安排明显不同。末期共有的4只为中国医药（600056）、南京高科（600064）、欧亚集团（600697）和中国平安（601318）；研报另有康力电梯、海康威视、老板电器和奥瑞金，当前资料下未通过本项目历史成长条件。第8章进一步区分已确认的规则差异、已定位的历史覆盖缺口与仍未解决的时点问题。

因此，本文将B定位为“保留原选股思想、明确月频近似的规则复现”，将E作为同规则历史延长。原研报36.74\%是其原样本与实施口径下的公布结果；本研究E的26.17\%另含2015年后的市场经历，不能据两者差额认定延展失败，也不能把B较高的历史年化解释为优于原模型。

% !TeX root = ../main.tex

\chapter{数据与变量}
\section{样本范围及可比性}
数据来自项目本地CSMAR合并表。股票收益保留市场类型1、4、16，估值截面沿用0、3、6开头的代码筛选。基础程序最后纳入2015年3月调仓，实际计收益至2015年6月；延展与四路径最后纳入2026年6月调仓，计收益至2026年9月。原始表覆盖终点与实际持有收益终点分别核查，不将未进入持仓计算的月份计入样本。

\begin{table}[htbp]
\centering
\caption{样本、字段与研究用途}
\renewcommand{\arraystretch}{1.18}
\small
\setlength{\tabcolsep}{4.5pt}
\begin{tabularx}{\textwidth}{@{}p{2.9cm}LL@{}}
\toprule
数据 & 主要字段或处理 & 用途 \\
\midrule
股票月度行情 & Mretwd；含现金红利再投资 & 组合月收益 \\
月末估值 & MarketValue；当月末总市值 & PE与规模 \\
净利润及EPS & B002000000；标准季度累计总利润 & 同比成长与TTM利润 \\
指数日行情 & 000001的Retindex；月内复利 & 上证综指基准 \\
冻结公告敏感性结果 & 2003-04至2024-12，261个月 & 附录信息时点诊断 \\
\bottomrule
\end{tabularx}
\end{table}

主比较分为两个层次。第一，147个月共同复现区间，基础复现、延展与四路径均使用相同月份。第二，282个月完整延展区间，以基础延展为同规则对照；基础复现不强行外推至2026年。年度表中2003年和2026年各只有9个月，其收益不当作完整自然年度收益。

\section{成长率、滚动利润与估值}
令NP为同一报告期的累计总净利润，成长率定义为：

\begin{equation}
g_{i,q}=\frac{NP_{i,q}}{NP_{i,q-4}}-1.
\label{eq:growth}
\end{equation}

去年同期利润非正时，成长率置为缺失，避免亏损扩大被机械识别为正增长。保留3、6、9、12月标准季末财报，不把1月1日的期初记录当作普通季度。非年报TTM利润由“上年全年利润$+$本期累计利润$-$去年同期累计利润”得到，年报本身直接使用全年利润。

\begin{equation}
PE_{i,t}=\frac{MV_{i,t}}{NP^{TTM}_{i,q}}.
\label{eq:pe}
\end{equation}

市值使用调仓当月值，利润与市值均有效且TTM非零时计算带符号PE。所有有限负PE参与升序排序并占用名次，最终入选要求PE为正。这严格沿用基础复现，并不等于先删除负PE、再在正PE内部取30\%。

\section{累计利润、TTM与同比的计算示例}
季度财报中的累计净利润包含年初至报告期末的利润，不能把一季度、半年度和前三季度累计数直接相加。举例而言，若上年全年净利润为100、今年半年度为70、上年同期半年度为50，则最近四个季度利润为100+70$-$50=120，累计利润同比为70/50$-$1=40\%。若同期总市值为2,400且单位相同，计算PE为20。这里的数字仅用于说明公式，不代表任一实际持仓；实际结果全部使用CSMAR原始字段。

同一条利润资料承担两种不同作用：同比衡量相对去年同期的成长，TTM为估值提供近一年盈利分母。二者不能互相替换。若上年同期为负，直接相除可能使扭亏、减亏或亏损扩大得到难以按通常正增长排序的数字，因此基础版将其同比置为缺失；缺失不表示增长为零，也不表示真实利润下降。历史“全为正”条件下，这条记录不能作为正成长通过。

市值除以总净利润与股价除以EPS具有相近的估值思想，但只有利润归属、股数定义及时间口径一致时才可严格互换。总净利润与归母净利润、期末总股本与EPS的加权平均股本均可能不同。本研究以基础程序口径为标准，不因另一字段提高收益就替换它；第8章的CSMAR归母利润、现成PE及行情表总市值对照，用于评估这些定义的敏感性。

\section{数据补齐解决了什么问题}
扩大时间覆盖能够增加可用月份和历史财报，避免延展结果因旧文件尾部缺失而停留在较早终点，但不自动解决每只股票的上市前财报覆盖、同比基数、财报修订和真实公告时间。对某一股票而言，“有20条报告记录”还不等于“有20个有效同比”：每一期同比还需要去年同期的利润基数，基数非正或缺失仍无法通过历史成长条件。

这也说明不能用“已补齐大部分时间区间”替代逐股诊断。本文同时检查输入版本、实际收益终点、财报连续性和同比有效数；原研报末期四只股票的缺口在第8章单独列明。当前证据只能说明本地资料按本规则会排除它们，不能推断这些公司当年真实经营不满足原作者条件。

\section{三维得分与财务信息边界}
质量指标为最近最多20条累计利润记录中净利润为正的比例，反映利润正值出现频率，不是ROE、现金流质量或单季利润率。动量为调仓月份之前12个月的累计收益，至少有8个月有效记录；规模为总市值，市值越小得分越高。三个指标在原有历史合格候选池内转换为方向一致的百分位后等权平均。

缺失质量按0、缺失动量按$-$1进入原有评分。四路径不改变此处理。财务指标直接按报告期映射到季度末，没有要求真实公告日已到达；使用最终会计值也未逐项恢复历史修订版本。因此，本文所有主表首先是既有基础口径下的回溯模拟。附录另以冻结的261个月公告敏感性结果说明这一限制。

% !TeX root = ../main.tex

\chapter{研究设计与路径定义}
\section{基础复现与同规则延展}
基础成熟期要求最近20季齐全、成长率全部为正，当前成长率居全截面有效成长率前30\%；估值在全截面有限PE中取最低30\%后要求为正，最后取交集。2003—2007年采用自首条成长记录起每期为正且无断档的初始替代要求，不预先限制成长率百分比，在最终交集内按成长率降序取最多15只；同值以股票代码升序决定先后。成熟期交集全部纳入，不设15只上限。

季度末名单从次月起使用三个月。因此，2007年12月初始期名单仍影响2008年第一季度收益，成熟期持仓从2008年4月开始。基础延展只延长计算终点，不改变上述规则。

\section{四条优化路径}
\begin{table}[htbp]
\centering
\caption{六个正式方案的规则关系}
\renewcommand{\arraystretch}{1.18}
\small
\setlength{\tabcolsep}{4.5pt}
\begin{tabularx}{\textwidth}{@{}p{2.5cm}*{4}{C}@{}}
\toprule
方案 & 成熟期历史门槛 & 成长率与PE排名 & 成熟期组合 & 动态仓位 \\
\midrule
B 基础复现 & 20/20为正 & 全截面分别排名 & 交集全部纳入 & 无 \\
E 基础延展 & 同B & 同B & 同B & 无 \\
P1 路径一 & 至少19/20为正 & 历史合格候选池内 & 两池合并打分，最多15只 & 无 \\
P2 路径二 & 同P1 & 全截面分别排名 & 同P1 & 无 \\
P3 路径三 & 同P1 & 同P1 & 名单同P1 & 原有风控 \\
P4 路径四 & 同P2 & 同P2 & 名单同P2 & 原有风控 \\
\bottomrule
\end{tabularx}
\end{table}

四路径均恢复正PE条件及基础版初始五年规则，初始名单不足15只时不按三因子补位。P1是正式的“成功选股法优化”，P2是第二条独立排名路径；P3、P4分别叠加在P1、P2上，不是在两种选股规则之间再次择优切换。

P1、P2成熟期的核心候选要求两项名次均达到30\%阈值，扩展候选使用50\%阈值。核心与补位合并后统一按三维得分截取最多15只，核心只是候选类别，没有优先保留权。因此，本文不写“原筛选核心全部保留”。E到P1同时包含历史门槛、排名、组合容量与三因子得分等变化，其差异只能作为整套组合规则的比较。

\begin{figure}[htbp]
\centering
\includegraphics[width=\linewidth]{fig01}
\caption{基础规则与四条优化路径的对照关系}
\par\smallskip
\begin{minipage}{0.97\linewidth}\small 注：箭头表示设计上的对照关系，不表示收益因果分解。资料来源：CSMAR，作者计算。\end{minipage}
\end{figure}

\section{排名母体的精确定义}
基础版分别用当季所有有限成长率股票建立成长率榜，用所有有限PE股票建立估值榜，记样本数为Ng、Np。先排序，再与历史条件取交集；两榜样本数不必相等。P2恢复该口径，核心阈值为$\lfloor 0.3N_g\rfloor$和$\lfloor 0.3N_p\rfloor$，扩展阈值相应为50\%。

P1则先筛出同时具备有限成长率、有限PE且历史合格的H只股票，再在其中建立两榜；两个30\%阈值都基于H。它的准确含义是“合格候选池内前30\%”，不能称为“全市场前30\%”。排名采用离散名次及向下取整，不是分位数插值。

以2024年12月末为例，Ng=3,974、Np=5,121、H=82，对应30\%截止名次分别为1,192、1,536和24。名次与分母同时变化，因此两种规则既可能互相纳入，也可能互相排除股票，不能预先认定全截面排名一定更严格或更宽松。P1与P2的历史合格池及三因子得分已逐条核对一致，以避免混入新的评分变化。

\section{排名分母如何改变实际入选}
以当前已核验的2024年末为例，全截面成长率榜有3,974只、PE榜有5,121只，历史合格且两指标有限的候选池仅82只。基础版及P2核心的两项截止名次分别为1,192和1,536；P1在82只内部的截止名次均为24。这里比较的是不同榜单上的位置，不能直接用“24小于1,192”推断某条规则更苛刻，因为单只股票的名次也随母体改变。

可以用假设例子理解：一只股票在合格池的成长率排第20、PE排第20，会通过P1两项核心阈值；若它在全市场成长率排第1,300，则不通过P2成长阈值。另一只股票在合格池成长率排第30，但在全市场排第1,000，且PE同时过关，则可能通过P2却不能通过P1。示例只说明逻辑可能性，实际组合变化仍以逐季名单核验为准。

负PE占位是另一个独立问题。若全市场有100条有限PE，其中25条为负，按基础规则先升序取最低30条，再要求PE为正，则正PE最多只剩5条，而不是正PE股票75只内部的最低30\%即22条。这不是说负PE股票可以买入，而是负PE参与构造比较门槛。把正PE筛选移动到排名之前，会改变阈值及名单，属于另一个实验；当前基础版与正式优化路径均不作这种隐含替换。

P1到P2的设计对照固定历史候选条件与得分，再改变成长和PE的排名母体，因而能更集中地观察排名实施效果。E到P1则还包含20/20变19/20、扩展池、三维得分和成熟期持仓上限变化；即使收益提高，也无法仅从这一组差额分离出某一个环节的贡献。

\section{动态仓位的实际机制}
风控首先用此前最多12个月底层未控净收益计算年化波动率，至少6个历史月后实施20\%目标波动率。仓位为目标波动除以历史波动并截于1，无杠杆；历史不足或波动不可用时维持满仓。

\begin{equation}
w_t=\min\left(1,\frac{0.20}{\hat{\sigma}_{t-1}}\right).
\label{eq:weight}
\end{equation}

另在每个新季度将参考净值重置为该季度起点净值。月初根据已完成的净值判断：若相对季度起点下跌达到25\%，该季度余下月份仓位置0；下一季度重置。代码未在季度内持续更新峰值，因此它是季度起点损失阈值，不能称为滚动最高净值回撤止损，也不能保证月内跌幅不超过25\%。

实际受控收益为底层已扣基础费用的净收益乘以仓位。现金收益记0，未另计动态仓位变化成本；基础费用也随仓位缩放，不是逐笔交易成本模型。P3和P4不改变持仓名单，减少的是风险暴露而非股票数量。

\section{收益实现与评价指标}
每月对有收益记录的目标持仓重新等权平均；当月没有有效收益则记0。每个有持仓季度首月固定扣0.13\%，不按真实换手率计算。该实现含月度再等权和缺失收益再平均，不是季度买入后权重自然漂移的严格模拟。

\begin{equation}
CAGR=\left[\prod_{t=1}^{T}(1+R_{p,t})\right]^{12/T}-1.
\label{eq:cagr}
\end{equation}

波动为月收益样本标准差乘$\sqrt{12}$；最大回撤包含初始净值1；夏普沿用程序零无风险利率口径；Calmar为CAGR除以最大回撤绝对值。所有跨方案主表统一月份。本文没有为多路径比较新增显著性检验，也不将较高夏普或CAGR差直接解释为独立alpha。历史19/20参数参考过后续表现，所有分段均为回溯描述，不称为未接触的样本外测试。

% !TeX root = ../main.tex

\chapter{复现一致性与全期表现}
\section{原复现区间的核对}
基础复现与延展在147个共同收益月份完全一致。四路径初始20个季度的256条股票—季度记录均与基础版逐股一致；P3、P4名单相同，但由于仓位管理，初始期收益不必与未控组合相同。

\begin{table}[htbp]
\centering
\caption{共同复现区间的绩效（2003-04至2015-06）}
\renewcommand{\arraystretch}{1.18}
\small
\setlength{\tabcolsep}{4.5pt}
\begin{tabularx}{\textwidth}{@{}p{2.55cm}*{4}{C}@{}}
\toprule
方案 & 月数 & 年化收益/\% & 波动/\% & 最大回撤/\% \\
\midrule
基础复现 B & 147 & 38.86 & 35.97 & -70.11 \\
基础延展 E & 147 & 38.86 & 35.97 & -70.11 \\
路径一 P1 & 147 & 41.73 & 35.85 & -68.87 \\
路径二 P2 & 147 & 43.33 & 36.14 & -64.15 \\
路径三 P3 & 147 & 30.61 & 23.48 & -32.74 \\
路径四 P4 & 147 & 30.92 & 24.39 & -34.26 \\
上证指数 & 147 & 9.09 & 28.93 & -70.97 \\
\bottomrule
\end{tabularx}
\end{table}

基础规则在共同区间的年化收益为38.86\%，最大回撤为$-$70.11\%。P2在该区间年化43.33\%，高于P1的41.73\%；后文可见，这一排序在完整延展区间逆转。因此，不能仅凭原复现终点选择所谓最优排名方式。本研究为月频近似，不把原研报的不同起止日期和日频结果当作逐日精确核对标准。

\begin{figure}[htbp]
\centering
\includegraphics[width=\linewidth]{fig02}
\caption{全部方案的共同区间净值}
\par\smallskip
\begin{minipage}{0.97\linewidth}\small 注：基础复现与延展共同月份完全重合，仅展示一条基础曲线；使用对数坐标。资料来源：CSMAR，作者计算。\end{minipage}
\end{figure}

\section{延展后的收益与风险}
\begin{table}[htbp]
\centering
\caption{完整延展区间的收益与风险（2003-04至2026-09）}
\renewcommand{\arraystretch}{1.18}
\small
\setlength{\tabcolsep}{4.5pt}
\begin{tabularx}{\textwidth}{@{}p{2.35cm}*{5}{C}@{}}
\toprule
方案 & 年化收益/\% & 波动/\% & 最大回撤/\% & 零利率夏普 & Calmar \\
\midrule
基础延展 E & 26.17 & 32.41 & -70.11 & 0.881 & 0.373 \\
路径一 P1 & 29.64 & 29.98 & -68.87 & 1.021 & 0.430 \\
路径二 P2 & 27.38 & 30.93 & -64.15 & 0.941 & 0.427 \\
路径三 P3 & 23.19 & 21.45 & -32.74 & 1.084 & 0.708 \\
路径四 P4 & 20.85 & 22.21 & -34.86 & 0.967 & 0.598 \\
上证指数 & 4.13 & 24.13 & -70.97 & 0.289 & 0.058 \\
\bottomrule
\end{tabularx}
\end{table}

282个月中，E年化26.17\%、波动32.41\%、最大回撤$-$70.11\%。P1年化29.64\%，较E增加3.48个百分点；波动降为29.98\%，但最大回撤仍达$-$68.87\%，仅收窄1.24个百分点。增加持股和提高夏普，并没有消除极端峰谷损失。

P2年化27.38\%、波动30.93\%、最大回撤$-$64.15\%。相较P1，它以2.26个百分点年化收益的下降，换取4.72个百分点的回撤收窄；波动反而上升0.95个百分点。故排名恢复没有同时改善收益、波动和回撤，也不能被概括为全面优于候选池内排名。

\begin{figure}[htbp]
\centering
\includegraphics[width=\linewidth]{fig03}
\caption{基础延展与两条选股优化路径的全期净值及回撤}
\par\smallskip
\begin{minipage}{0.97\linewidth}\small 注：上图为对数净值，下图为从初始净值1计算的月末回撤；同一区间282个月。资料来源：CSMAR，作者计算。\end{minipage}
\end{figure}

\section{仓位管理的成对影响}
P3将P1波动由29.98\%降至21.45\%，最大回撤由$-$68.87\%收窄至$-$32.74\%；代价是年化从29.64\%降至23.19\%。P4将P2波动由30.93\%降至22.21\%，最大回撤由$-$64.15\%收窄至$-$34.86\%，年化从27.38\%降至20.85\%。两组均体现风险暴露下降与上涨参与下降并存。

P3零利率夏普1.084，高于P1的1.021；P4为0.967，高于P2的0.941。风控改善了所设夏普，但没有提高全期绝对收益。P3在本次完整区间同时高于P4的收益、低于其波动和回撤；这只是指定规则及样本中的排序，并非对未来或所有风险偏好的支配结论。

\begin{figure}[htbp]
\centering
\includegraphics[width=\linewidth]{fig04}
\caption{两条风控路径的全期净值及回撤}
\par\smallskip
\begin{minipage}{0.97\linewidth}\small 注：基础延展不额外叠加风控，作为相同信息口径的原规则对照。资料来源：CSMAR，作者计算。\end{minipage}
\end{figure}

\begin{figure}[htbp]
\centering
\includegraphics[width=\linewidth]{fig05}
\caption{各正式路径的年化收益与最大回撤}
\par\smallskip
\begin{minipage}{0.97\linewidth}\small 注：斜线柱为动态仓位路径；最大回撤越接近零表示峰谷损失越小。资料来源：CSMAR，作者计算。\end{minipage}
\end{figure}

\section{如何理解收益提高与风险降低}
全期年化收益回答的是持有整个样本的复利速度，年化波动描述月收益离散程度，最大回撤描述历史峰值到后续谷值的最大损失。三者不是同一风险收益评价。P1相对E同时提高年化、降低波动，但最大回撤只改善1.24个百分点，说明常态收益分布有所改善仍可与极端损失集中出现并存。不能把波动下降或持股增加直接表述为历史大跌已被化解。

P1与P2的比较同样需要多指标阅读：P2全期收益较低、最大回撤较小，月收益波动却略高。因此，“风险更低”必须明确指哪项风险。若研究目标是提高绝对复利，当前全样本结果支持P1居前；若关注原始排名解释性，P2提供更接近基础版的母体定义；若主要控制本样本峰谷损失，则需要连同P3、P4的收益代价评价，不能把三种目标合并为单一“最好”。

累计财富差异会被长期复利放大。P1与P2的年化只相差约2.26个百分点，经过23.5年，终值比例仍可能有明显差距；这种放大不意味着每年都领先，也不意味着收益差具有统计显著性。后续年度和分段表负责检验领先是否集中于少数时期，而非根据全期终值给出稳定性结论。

\section{研报原始时期内的优化效果检验}
\subsection{固定截止范围与比较窗口}
此前共同147个月的结果以2015年6月结束，而原研报公布的结果至2015年7月31日。为检验优化是否在原始时期内同样改善，本节从CSMAR逐股月收益独立重算，统一保留2015年5月31日前的选股信号；按本项目季度末规则，最后信号为3月31日，共49次。末组持仓由4—6月延持至7月，不加入6月新信号，不在没有新交易的7月另扣一次季度费用。因此，本节没有用原资料截止后才形成的新季度信号填充最后一个月。

评价起止分为三种：147个月为2003年4月至2015年6月，与既有基础版对应；148个月保留相同起点并延长末组至7月；151个月在148个月前加入2003年1—3月建仓前现金。现金仅用于起始分母敏感性，不作为已确认的原研报建仓规则，也不给风控增加虚构的预热收益。上证指数在1—3月使用实际CSMAR收益，不跟随策略置为现金。

风控严格沿用“选股信号变化时重置参考净值”的原机制。延持7月没有新信号，故仍使用末组参考值；历史波动窗口继续按此前收益更新。成熟期从2008年4月至2015年7月，共88个月，控制路径继承先前状态，不在2008年重新启动。月末回撤在每个评价窗口包含初始净值1。

\begin{table}[htbp]
\centering
\caption{原始时期不同起止口径的年化收益}
\renewcommand{\arraystretch}{1.18}
\small
\setlength{\tabcolsep}{4.5pt}
\begin{tabularx}{\textwidth}{@{}p{2.55cm}*{4}{C}@{}}
\toprule
方案 & 147月/\% & 148月/\% & 151月/\% & 151月相对B/百分点 \\
\midrule
基础复现B & 38.86 & 37.31 & 36.44 & +0.00 \\
路径一P1 & 41.73 & 39.70 & 38.77 & +2.33 \\
路径二P2 & 43.33 & 41.95 & 40.97 & +4.52 \\
路径三P3 & 30.61 & 29.08 & 28.43 & -8.01 \\
路径四P4 & 30.92 & 30.14 & 29.46 & -6.98 \\
上证指数 & 9.09 & 7.67 & 8.42 & -28.02 \\
\bottomrule
\end{tabularx}
\par\smallskip
\begin{minipage}{\textwidth}\small 注：147月、148月、151月并非三次独立投资试验，而是同一原始时期的起止敏感性；相对B为同一151个月内的CAGR差。申万公布的36.74\%仍具有不同公告日调仓及日频定义，本表不把它当作已经同口径的基准。\end{minipage}
\end{table}

\subsection{收益改善及其风险代价}
三种窗口均出现P1、P2年化高于B，P3、P4年化低于B的方向。在151个月口径下，B年化36.44\%，P1为38.77\%、P2为40.97\%，分别提高2.33和4.52个百分点。将终点从6月推至7月后，这一选股改善方向仍保留，而不是仅在某一个截止月份成立。P2在原始时期高于P1，亦与后文完整282个月中P1领先的排序不同，进一步说明规则的相对表现依赖市场阶段。

P1、P2的最大回撤分别为$-$68.87\%、$-$64.15\%，相对B的$-$70.11\%收窄1.24、5.96个百分点。但年化波动与B接近或略高，所以“收益和最大回撤改善”不等于全部风险指标同步改善。尤其P1仍有接近69\%的历史峰谷损失，不能因年化增加就忽略损失承受要求。

P3、P4年化降至28.43\%、29.46\%，较B下降8.01、6.98个百分点；回撤则收窄至$-$32.74\%、$-$34.26\%，零利率夏普分别为1.187、1.195，高于B的1.055。这支持风控在原始时期降低风险暴露的描述，却不支持其提高全原始区间绝对收益。此处P4的$-$34.26\%与完整延展的$-$34.86\%来自不同终点，应分别阅读。

\begin{table}[htbp]
\centering
\caption{原始时期风险及成熟期的补充检验}
\renewcommand{\arraystretch}{1.18}
\small
\setlength{\tabcolsep}{4.5pt}
\begin{tabularx}{\textwidth}{@{}p{2.35cm}*{5}{C}@{}}
\toprule
方案 & 151月回撤/\% & 151月波动/\% & 151月夏普 & 成熟期年化/\% & 成熟期回撤/\% \\
\midrule
基础复现B & -70.11 & 35.74 & 1.055 & 18.45 & -62.72 \\
路径一P1 & -68.87 & 35.74 & 1.102 & 21.94 & -60.64 \\
路径二P2 & -64.15 & 35.86 & 1.143 & 25.26 & -54.33 \\
路径三P3 & -32.74 & 23.53 & 1.187 & 19.10 & -26.15 \\
路径四P4 & -34.26 & 24.18 & 1.195 & 20.75 & -34.26 \\
上证指数 & -70.97 & 28.99 & 0.425 & 1.08 & -53.19 \\
\bottomrule
\end{tabularx}
\par\smallskip
\begin{minipage}{\textwidth}\small 注：成熟期88个月为2008-04至2015-07；分段回撤按段初净值1重算，不能替代跨阶段回撤。夏普沿用零无风险利率口径，收益、回撤与波动为百分数。\end{minipage}
\end{table}

\subsection{排除共同初始阶段后如何判断}
成熟期88个月中，B年化18.45\%，P1、P2为21.94\%、25.26\%，仍分别提高3.49、6.81个百分点。因此，两条选股改进的优势并非只由完全相同的初始五年名单产生。P3、P4在该段年化19.10\%、20.75\%，也分别高于B约0.65、2.30个百分点，但仍低于各自底层P1、P2。

这说明“风控在整个原始时期降低收益”与“成熟期风控路径收益高于基础版”可以同时成立：前者与原始B比较并包括初期高涨，后者只评价成熟期，且底层选股也发生变化。不能据成熟期P3高于B，就把收益增加单独归于风控；识别仓位作用仍应比较P1与P3、P2与P4。初始高涨参与减少与阶段排序共同提供解释，但没有构成各机制可加总的因果分解。

本实验24行指标覆盖6个方案和4个窗口，R逐股重算后与原147个月逐月一致；Python再次独立连接CSMAR收益、重建148个月控制仓位并复算指标。其证据支持指定月频口径下的改善方向和计算一致性，尚未解决真实公告时点、原历史覆盖及成交约束，也不是事先冻结的样本外验证。完整实验程序和逐月结果独立保存在“研报原始区间优化验证\_20261107”。

\begin{figure}[htbp]
\centering
\includegraphics[width=\linewidth]{fig06}
\caption{原始时期各路径的收益与最大回撤}
\par\smallskip
\begin{minipage}{0.97\linewidth}\small 注：151个月；建仓前现金及末组延持仅为明示敏感性，不是申万精确日频重建。资料来源：CSMAR，作者计算。\end{minipage}
\end{figure}

\begin{figure}[htbp]
\centering
\includegraphics[width=\linewidth]{fig07}
\caption{原始时期的月度净值与回撤过程}
\par\smallskip
\begin{minipage}{0.97\linewidth}\small 注：实际投资期148个月；统一2015年5月前信号，不使用6月新信号。资料来源：CSMAR，作者计算。\end{minipage}
\end{figure}

% !TeX root = ../main.tex

\chapter{历史分段与收益持续性}
\section{初始期、成熟期与追加区间}
\begin{table}[htbp]
\centering
\caption{不同历史区间的年化收益（\%）}
\renewcommand{\arraystretch}{1.18}
\small
\setlength{\tabcolsep}{3pt}
\begin{tabularx}{\textwidth}{@{}p{3.15cm}*{6}{C}@{}}
\toprule
区间 & E & P1 & P2 & P3 & P4 & 上证 \\
\midrule
2003-04至2008-03 & 70.52 & 70.52 & 70.52 & 45.25 & 45.25 & 18.11 \\
2008-04至2015-06 & 20.53 & 24.76 & 27.14 & 21.38 & 21.86 & 3.27 \\
2015-07至2024-12 & 9.95 & 15.10 & 12.19 & 13.26 & 10.65 & -2.62 \\
2025-01至2026-09 & 36.03 & 32.53 & 11.18 & 29.07 & 11.48 & 8.12 \\
\bottomrule
\end{tabularx}
\end{table}

初始期60个月，E、P1、P2均年化70.52\%，反映三者完全相同的早期选股和特定市场路径。其高复利对全期财富积累具有较大影响，不能把全期约30\%的年化描述为策略进入成熟市场后一直保持的增长速度。风控路径在同段降为45.25\%，也显示缩仓会放弃部分上涨收益。

成熟期前段P2年化27.14\%，高于P1的24.76\%；后续2015年7月至2024年12月则P1为15.10\%、P2为12.19\%，排序逆转。追加21个月中，E年化36.03\%，高于P1的32.53\%、P3的29.07\%、P4的11.48\%和P2的11.18\%。该区间较短，年化换算对短期路径敏感，但已足以反驳“四条优化在每一时期都优于基础延展”的表述。

分段回撤在各段初始净值1上重新计算，不能直接替代跨阶段的全期回撤。历史分段、参数选择和研究者已见数据的关系也限制预测解释，本表不承担严格样本外验证功能。

\section{年度表现与滚动窗口}
\begin{figure}[htbp]
\centering
\includegraphics[width=\linewidth]{fig08}
\caption{年度复利收益的跨方案比较}
\par\smallskip
\begin{minipage}{0.97\linewidth}\small 注：数字为实际年度复利百分比，色阶仅为便于阅读截于$\pm$60\%；2003年与2026年各含9个月。资料来源：CSMAR，作者计算。\end{minipage}
\end{figure}

年度热图显示，同一时期的市场环境对各路径有共同影响。风控路径并不逐年增加收益，选股路径的相对排序也随年份变化。完整年度收益详列附录A，避免只呈现有利年份。2026年结果必须与“前9个月”的统计长度一起解读。

\begin{figure}[htbp]
\centering
\includegraphics[width=\linewidth]{fig09}
\caption{滚动五年年化收益}
\par\smallskip
\begin{minipage}{0.97\linewidth}\small 注：每个点使用此前60个月；窗口相互重叠，不能作为彼此独立的观察或胜率检验样本。资料来源：CSMAR，作者计算。\end{minipage}
\end{figure}

滚动曲线用于显示收益水平随历史终点变化，而不是重复筛选最优参数。高全期CAGR与较低后段收益可以同时成立；以固定全期均值替代各阶段经验，会掩盖早期复利与后期回报的差别。

\section{初始高收益与后续延展应当怎样连接}
初始期60个月的高收益来自三条未控规则完全相同的选股记录，因此不能作为成熟期优化的证据。它会共同抬升E、P1和P2的完整区间年化，使一个成熟期表现一般的路径仍呈现较高全期数值。为此，成熟期合计还应独立阅读：E、P1、P2年化分别为16.30\%、20.39\%、17.73\%，均低于各自包含初始期的完整年化。

这不是把早期收益删除，而是区分完整投资路径与规则改变后阶段。2007年12月产生的初始规则名单持续影响2008年1至3月，所以分界必须按实际持有月份放在2008年4月，不能仅按自然年2008年1月切开。此处的月份连接与程序的“季度末信号、后三个月持有”一致。

2015年7月至2024年12月，P1年化15.10\%高于P2的12.19\%；追加21个月中，E以36.03\%领先，P1为32.53\%、P2仅11.18\%。原区间与延展区间的排名逆转有实际数据依据，因此本研究同时保留两条路径，不按某一个终点重新决定唯一正确策略。追加区间亦是已知历史的一部分，不能被包装为事前冻结后的独立预测验证。

\section{年度优势与回撤恢复的补充观察}
附录D进一步列出完整自然年度的正收益数、领先基础延展的年度数以及最差年度；只纳入2004—2025年各12个月的完整年度。2003年和2026年不混入年度计数，避免长度差异。年度次数只是描述，对照中的年份也共享市场环境，不能解释为未来获胜概率或独立试验。

同时，最大回撤幅度之外还应关注从峰值到谷值、再到首次恢复峰值的时间。一条路径即使最大损失较小，也可能较长时间未创新高。附录D以每月末净值计算相关日期和月数，初始净值纳入峰值搜索；这比单一回撤百分数多提供时间维度，但仍无法观察月内低点或投资者实际申赎路径。

% !TeX root = ../main.tex

\chapter{排名母体与组合构成机制}
\section{排名变化的描述性幅度}
\begin{table}[htbp]
\centering
\caption{成对调整的描述性变化}
\renewcommand{\arraystretch}{1.18}
\small
\setlength{\tabcolsep}{4.5pt}
\begin{tabularx}{\textwidth}{@{}p{1.45cm}*{4}{C}@{}}
\toprule
对照 & 年化变化/百分点 & 波动变化/百分点 & 回撤变化/百分点 & 夏普变化 \\
\midrule
E$\to$P1 & 3.48 & -2.43 & 1.24 & 0.140 \\
P1$\to$P2 & -2.26 & 0.95 & 4.72 & -0.080 \\
P1$\to$P3 & -6.46 & -8.53 & 36.13 & 0.063 \\
P2$\to$P4 & -6.53 & -8.72 & 29.29 & 0.026 \\
\bottomrule
\end{tabularx}
\end{table}

表10中的E到P1是整套组合规则变化，P1到P2才对应固定历史池和得分后的排名母体调整；P1到P3、P2到P4对应同一底层组合的仓位管理。回撤变化取“新值$-$旧值”，正值表示负回撤收窄，不能把它误读为风险增加。收益、波动和回撤变化用百分点，不是相对百分比。

\begin{figure}[htbp]
\centering
\includegraphics[width=\linewidth]{fig10}
\caption{收益—波动位置与成对风控调整}
\par\smallskip
\begin{minipage}{0.97\linewidth}\small 注：箭头分别连接P1到P3、P2到P4；该图不表示均值—方差有效前沿。资料来源：CSMAR，作者计算。\end{minipage}
\end{figure}

P1与P2的最终名单在成熟期74个调仓日均发生变化。两种排名既改变核心范围，也改变扩展候选，最后通过统一三因子截取改变持仓。全截面排名不能独立于最终组合规则解释，其收益差是相同程序其余环节下的实施效果。

\begin{figure}[htbp]
\centering
\includegraphics[width=\linewidth]{fig11}
\caption{排名调整的相对净值与持仓集合相似度}
\par\smallskip
\begin{minipage}{0.97\linewidth}\small 注：相对净值为P2净值除以P1净值；Jaccard为两名单交集数除以并集数。该相似度不是成交换手率。资料来源：CSMAR，作者计算。\end{minipage}
\end{figure}

\section{持股数量、缺失收益与当前成长率}
\begin{table}[htbp]
\centering
\caption{持仓覆盖与当前成长率诊断}
\renewcommand{\arraystretch}{1.18}
\small
\setlength{\tabcolsep}{4.5pt}
\begin{tabularx}{\textwidth}{@{}p{1.55cm}*{5}{C}@{}}
\toprule
方案 & 平均持股/只 & 空仓信号/次 & 记录数/条 & 收益缺失月/个 & 当期非正成长/条 \\
\midrule
E & 5.96 & 8 & 560 & 5 & 0 \\
P1 & 14.27 & 0 & 1341 & 16 & 0 \\
P2 & 14.53 & 0 & 1366 & 36 & 17 \\
\bottomrule
\end{tabularx}
\end{table}

E平均持股5.96只、8个空仓信号；P1、P2分别为14.27和14.53只，均无空仓，但这不等于每季固定15只，初始期仍可能不足。E成熟期交集全部纳入，最大持股18只；P1、P2成熟期由统一得分截于15只。

目标持股多于有效月收益的月份，E有5个，P1有16个，P2有36个。所有路径沿用相同缺失收益再平均方式，但扩展候选改变了缺失事件发生频率，因而不能认为缺失收益对路径比较必然无影响。未有收益记录的股票不应被解释为已经以合理价格卖出；停退市真实损失尚未还原。

P2有17条股票—季度记录的当前成长率非正。这与20季至少19季为正的历史条件并不矛盾，也与正PE条件不同。全市场前50\%在盈利普遍下降时可以包含负成长企业，因此本路径保留的是较宽的历史持续性约束，而非要求每一只持仓当季都在增长。P1本次未出现这类持仓，亦不能据此把其程序门槛改写为“当前成长必须为正”。

\section{候选核心与最终核心}
\begin{table}[htbp]
\centering
\caption{成熟期核心候选与最终持仓的区别}
\renewcommand{\arraystretch}{1.18}
\small
\setlength{\tabcolsep}{4.5pt}
\begin{tabularx}{\textwidth}{@{}p{1.55cm}*{4}{C}@{}}
\toprule
路径 & 候选核心/条 & 最终核心/条 & 候选核心保留率/\% & 持仓核心占比/\% \\
\midrule
P1 & 439 & 309 & 70.39 & 28.48 \\
P2 & 552 & 137 & 24.82 & 12.34 \\
\bottomrule
\end{tabularx}
\end{table}

P2候选核心由439条增至552条，但最终保留核心由309条降至137条。候选核心条数跨季度累计，不是独立公司数；候选核心保留率以候选核心为分母，最终持仓核心占比以全部最终成熟期持仓记录为分母，两个比例不可混用。

这一变化与补位池从1,097条扩大到3,310条、最终又统一按得分截取的机制一致。扩展候选可以挤出核心，较大的核心池不保证最终核心比例提高。当前正式路径保留了用户指定的组合实现，因此本文不把结果称为“在不删任何核心股票的前提下改进”。

\begin{figure}[htbp]
\centering
\includegraphics[width=\linewidth]{fig12}
\caption{目标持股与实际持仓核心比例}
\par\smallskip
\begin{minipage}{0.97\linewidth}\small 注：核心比例仅用最终持仓标记计算；成熟期从2008年第一笔成熟期信号开始，基础原规则没有补位分层。资料来源：CSMAR，作者计算。\end{minipage}
\end{figure}

\section{为什么补位与核心保留必须分开讨论}
扩展候选的经济动机是缓解严格历史条件和两项30\%交集造成的持仓过少：在仍保留持续成长、相对估值和正PE约束的前提下，从较宽范围寻找替代股票。质量、动量与规模用于统一排序，它们只是对候选之间作选择，不能替代最初的成长与估值筛选。因此，这些改进仍保留原选股思想，但已经不是原交集的简单复制。

代码目前把候选核心与扩展候选合并，再统一按得分取最多15只。只要扩展候选得分更高，就可能排在核心前面；“核心”这个标签本身没有优先保留权。P2核心候选增加而最终核心减少，正好提醒研究者：扩大一个上游集合不保证扩大下游结果，最终截取规则同样重要。若日后要研究核心优先，应独立保存并另作对照，不能只修改报告文字来宣称已经实现。

持股数量增加也不等于完成风险分散证明。多只股票可能共享行业、规模或盈利周期暴露，现有结果只直接说明名单容量和部分风险统计的变化。本文没有据此宣称行业中性或独立风险已消除；原研报的行业比例不能直接代替当前组合的行业证据。

% !TeX root = ../main.tex

\chapter{动态仓位的作用与边界}
\section{风控是否实际触发止损}
\begin{table}[htbp]
\centering
\caption{风控实际仓位与月度风险诊断}
\renewcommand{\arraystretch}{1.18}
\small
\setlength{\tabcolsep}{4.5pt}
\begin{tabularx}{\textwidth}{@{}p{1.55cm}*{4}{C}@{}}
\toprule
路径 & 平均仓位/\% & 零仓位月/个 & 最差单月/\% & 上证下跌月均收益/\% \\
\midrule
P1 & 100.00 & 0 & -30.10 & -3.20 \\
P2 & 100.00 & 0 & -28.82 & -3.28 \\
P3 & 80.40 & 0 & -14.02 & -2.29 \\
P4 & 78.62 & 0 & -13.45 & -2.37 \\
\bottomrule
\end{tabularx}
\end{table}

P3平均仓位80.40\%，P4为78.62\%，两者零仓位月份均为0。按代码机制，若季度止损清仓实际生效，应出现零仓位月份；本样本没有该现象。因此，不能把最大回撤下降解释为25\%季度止损已经成功阻止损失，已观察到的风控作用主要体现为波动目标下的持续缩仓。对止损阈值本身的有效性仍缺少实际触发证据。

\begin{figure}[htbp]
\centering
\includegraphics[width=\linewidth]{fig13}
\caption{动态风控的月度目标仓位及年度平均仓位}
\par\smallskip
\begin{minipage}{0.97\linewidth}\small 注：仓位只使用月初已知历史收益；未加杠杆，现金利息为零。零仓位月数为零。资料来源：CSMAR，作者计算。\end{minipage}
\end{figure}

\section{收益代价与月度尾部风险}
P1最差单月为$-$30.10\%，P3为$-$14.02\%；P2为$-$28.82\%，P4为$-$13.45\%。在上证下跌月份，两组受控收益平均损失也较未控组合收窄。它们支持“缩仓降低本样本月度风险暴露”的描述，仍不能推出月内损失受到严格上限保护。

两条风控路径全期年化均比底层下降约6.5个百分点，说明风险降低不是无成本增加收益。实际部署还会受到现金利息、动态交易费用、持仓估值与可成交性的影响；本次没有引入这些机制，不能把当前缩放净收益当作完整账户回测。

\section{风控收益为什么会下降}
动态仓位在多数高波动月份减少暴露，既会缩小负收益，也会按相同比例缩小随后上涨。波动率依赖过去最多12个月的收益，市场快速反弹时，窗口内仍可能保留此前的高波动记录，仓位不能立刻恢复到1。因此，受控组合与底层组合名单相同，长期财富却可能明显不同；这是仓位机制的结果，不能把全部收益下降说成选股能力退化。

P3与P4的平均仓位分别为80.40\%和78.62\%，但不能把底层CAGR简单乘以平均仓位来预测受控CAGR。每个月的仓位与当月收益组合后再复利，仓位发生在哪些上涨或下跌月份，比单一平均值更重要。实际收益还沿用“底层净收益乘仓位”，基础费用一并缩放，现金利息为零；因此它是当前规则下的收益缩放模拟，而非完整交易账户。

风控改善也不在所有阶段同样显著。例如，2015年7月至2024年12月，P1与P3分段最大回撤均约29.42\%，P2与P4均约34.12\%。这说明完整区间回撤大幅收窄，不能直接推广为每一段下跌都得到等比例保护。附录D分别记录各路径最深回撤所在时段，有助于避免只比较某个历史百分数。

% !TeX root = ../main.tex

\chapter{与申万研报的差异归因及计算合理性证据}
\section{收益比较必须先对齐区间}
原研报[1]第6页报告截至2015年7月31日的2938个交易日、累计收益4176.87\%、年化36.74\%；本研究E的26.17\%覆盖到2026年9月。二者时期不同，不能把差额全部解释为复现错误。当前B在147个月的年化为38.86\%，实际高于36.74\%，但也不因此认定复现完全正确。

\begin{table}[htbp]
\centering
\caption{原研报与本研究收益区间的区别}
\renewcommand{\arraystretch}{1.18}
\small
\setlength{\tabcolsep}{4.5pt}
\begin{tabularx}{\textwidth}{@{}p{2.9cm}*{3}{C}@{}}
\toprule
对象与口径 & 评价月份 & 年化收益/\% & 累计收益/\% \\
\midrule
申万原文 & 原文称2003年至2015-07 & 36.74 & 4176.87 \\
既定基础复现 & 2003-04至2015-06 & 38.86 & 5481.49 \\
末组延持至7月 & 2003-04至2015-07 & 37.31 & 4891.07 \\
年初现金并延持7月 & 2003-01至2015-07 & 36.44 & 4891.07 \\
原规则完整延展 & 2003-04至2026-09 & 26.17 & 23457.11 \\
\bottomrule
\end{tabularx}
\par\smallskip
\begin{minipage}{\textwidth}\small 注：前两行不是完全同日同频比较。末组延持保留原截止前49个信号，不加入2015年6月信号，仅延长末组至7月；年初现金把1—3月未建仓收益置0。两行是起止月份敏感性，不是申万日频重建。\end{minipage}
\end{table}

终点和年化分母即可将38.86\%改变为37.31\%或36.44\%，但延持后的累计收益4891.07\%仍与4176.87\%不同。原文累计财富与年化反推约12.003年，而未披露全部年化日数约定；本文不据此认定原报告计算错误，也不把图中首个刻度当作已认证的建仓日。

\section{其他CSMAR字段与实施规则的独立对照}
本轮以现有CSMAR数据独立构造17组明确口径，在指定区间保存73行绩效。它们包括基准、单项敏感性、实施联合对照和错误字段反例，不是17个独立投资模型，也未进行组合参数寻优。表15列示主要对照，其余结果在配套分析文档保存。

\begin{table}[htbp]
\centering
\caption{统一区间后的主要口径敏感性}
\renewcommand{\arraystretch}{1.18}
\small
\setlength{\tabcolsep}{4.5pt}
\begin{tabularx}{\textwidth}{@{}p{3.35cm}*{2}{C}@{}}
\toprule
实现对照 & 147个月年化/\% & 282个月年化/\% \\
\midrule
既定月频基准 & 38.86 & 26.17 \\
行情表总市值 & 38.17 & 24.39 \\
归母利润计算PE & 44.63 & 26.74 \\
数据库同比独立历史 & 41.28 & 27.52 \\
单季同比成长 & 27.96 & 16.25 \\
五个完整年度历史 & 36.22 & 24.17 \\
财务固定滞后一季 & 18.62 & 13.39 \\
季度内权重漂移 & 38.73 & 25.82 \\
\bottomrule
\end{tabularx}
\end{table}

行情表总市值替代后原区间38.17\%，与基准较接近；归母利润计算PE为44.63\%，并未自然缩小差异。现成PE1/PE2在亏损时为空，需与同样采用正PE内排名的对照比较，不能将负PE占位变化混同于利润定义效应。

CSMAR字段字典明确F081001B是单季度环比、F081002B才是同比。原程序自行计算累计净利润同比，不存在误用字段A的问题。本轮另从04表自身建立正确同比历史，而非先丢掉03表没有的历史报告，年化为41.28\%；仍未解决末期4只股票缺失。

固定滞后一季使原区间年化降为18.62\%，表明财务期安排会改变表现，但一季滞后不等于原研报发布日算法。缺失收益按账面不变处理为38.74\%，再允许季度内权重漂移为38.73\%，本组机械收益实现影响较小。五个完整年度历史得到36.22\%，却在2015年3月选出20只、仅重合4只；把费用改到0.50\%也会接近原年化，却违背原费率。这些反例说明不能以“接近36.74\%”选择算法。

\begin{figure}[htbp]
\centering
\includegraphics[width=\linewidth]{fig14}
\caption{原区间单项口径验证}
\par\smallskip
\begin{minipage}{0.97\linewidth}\small 注：各柱均为相同147个月；较接近原报告数值不代表持仓一致。资料来源：CSMAR，作者计算。\end{minipage}
\end{figure}

\section{末期持仓差异可以定位到历史资料}
原研报第12页给出2015年5月25日调仓、7月31日持有8只。当前2015年3月信号为9只，重合4只；改用5月末市值配一季报可选出8只，仍仅重合4只。因此，数量相同不能替代名单核对。

\begin{table}[htbp]
\centering
\caption{未匹配四只股票的历史成长条件}
\renewcommand{\arraystretch}{1.18}
\small
\setlength{\tabcolsep}{4.5pt}
\begin{tabularx}{\textwidth}{@{}p{2.45cm}*{3}{C}@{}}
\toprule
股票 & 五年窗口财报数 & 正成长率数 & 缺失同比/期 \\
\midrule
康力电梯 & 20 & 17 & 3 \\
海康威视 & 20 & 16 & 4 \\
老板电器 & 18 & 14 & 4 \\
奥瑞金 & 10 & 6 & 4 \\
\bottomrule
\end{tabularx}
\end{table}

四只股票当期成长和PE名次均可通过原条件，阻碍来自历史资料：康力电梯、海康威视分别有3、4期缺少去年同期利润；老板电器、奥瑞金既缺报告期，也有4期缺同比基数。这些记录不应被说成实际负成长。04表正确同比补充后仍有缺失；05公告利润也未覆盖足够的上市前季度。

\begin{figure}[htbp]
\centering
\includegraphics[width=\linewidth]{fig15}
\caption{研报最后8只股票的条件核查}
\par\smallskip
\begin{minipage}{0.97\linewidth}\small 注：条件按当前CSMAR资料与2015年3月信号计算，不是对2015年5月25日日频价格的精确复原。资料来源：CSMAR，作者计算。\end{minipage}
\end{figure}

证据支持“当前历史覆盖会排除研报中的4只”，但不能进一步证明原作者忽略缺失、采用年频历史或拥有何种不同数据库。合理的下一步是核查CSMAR内上市前比较期财务及原文的数据说明，不是补零、豁免条件或强行加入这4只。

\section{调仓算法差异与计算证据的边界}
原研报第6页使用全市场最晚发布日与报告期后90个交易日孰早；本研究使用季度末报告期映射。这是已确认的复现近似，可能同时改变信号、成交价格和持有月份。按当前05表及所查代码范围计算，2015年一季报最晚有效公告为5月4日，原报告调仓为5月25日；共同日期仍不一致，不能直接硬填成原日期。

本轮没有2003—2015年个股日频成交重建，也未恢复当时财务修订版本。故本文证明的是既定月频定义下的计算可复核性，而非与原日频策略逐日等价。

\begin{table}[htbp]
\centering
\caption{程序与计算合理性的证据链及适用范围}
\renewcommand{\arraystretch}{1.18}
\small
\setlength{\tabcolsep}{4.5pt}
\begin{tabularx}{\textwidth}{@{}p{2.35cm}LL@{}}
\toprule
核验层次 & 实际证据 & 可以支持的结论 \\
\midrule
输入与变量 & CSMAR字段字典、单位、主键及季度累计利润核查 & 指标定义及连接具有可追溯依据 \\
独立基准重建 & 独立Python逐股重建E的560条持仓；282个月最大差约$5\times10^{-16}$ & 当前基准R实现与独立算法一致 \\
公式交叉复算 & R直接复算17组共同147个月收益与回撤，误差小于$10^{-10}$ & 实验汇总不是由单一程序自证 \\
四路径机制核验 & 初始名单、排名、三因子得分及过去收益仓位逐项检查 & 改进路径按已披露规则执行 \\
时点与资料边界 & 原调仓日及4只股票历史缺口尚未重建一致 & 不能证明申万日频等价或交易可实现性 \\
\bottomrule
\end{tabularx}
\end{table}

从方法上，季度财务构造成长与估值、月收益评价持仓，是与本项目月频近似目标相适应的设计；市值、利润、同比基数和收益单位均有定义，分母非正及缺失的处理也已公开。其合理性是“在明示目标和假设下可解释”，不是“适用于所有交易问题”。

从实现上，独立重建与R结果在持仓及浮点精度内一致，另有不同语言的公式复算和风控序列核查，足以支持“目前未发现既定规则实现中的重大算术或连接不一致”。独立程序也可能共享同一原始数据偏差，因此仍不能推出数据绝对正确、历史可交易或存在稳健超额收益。程序一致性与投资有效性必须分别评价。

\section{差异证据应按确定程度组织}
第一类是可直接确认的定义差异，包括本研究固定季度末、原研报发布日规则，以及实际收益终点分别为6月与7月。它们已由原文和源程序对应确认，不需要依靠收益接近程度来推测。第二类是当前数据支持的排除机制：末期四只股票存在历史同比或报告期缺口，按照本项目20/20条件必然无法入选。它说明本地实现为什么排除它们，但没有证明原作者采用了相同的历史覆盖或同一种缺失处理。

第三类是敏感性结果。归母利润、现成PE、财务滞后或完整年度历史等替代定义改变收益，说明输出对口径具有依赖性；它们不能反向认证哪一个定义就是原报告真正使用的版本。尤其是五个完整年度历史的年化36.22\%接近36.74\%，但末期名单仍明显不同，构成“只对收益率”的明确反例。

第四类仍待查明，包括原始数据库历史版本、精确首笔建仓日、年化日数约定及逐笔交易实现。当前证据不能给出各原因对最终差额的可加总百分比分解：某一个变化还会改变排名、名单及后续复利，多个变化并非线性叠加。本文因此使用定义对照、单项敏感性及逐股诊断共同解释，而不把CAGR差机械分摊成若干所谓贡献率。

\section{如何判断“计算合理”而不陷入自我证明}
输入层需要明确字段、单位、主键和财报期间，公式层需要说明利润同比和TTM如何构造，实施层需要检查排名、历史条件及名单，结果层需要从逐月收益复算复利和回撤。只有最终收益接近原文，无法替代前面各层检查。反过来，程序内部一致也不能认证数据供应商每条记录或公告时点都正确。

本研究独立Python从CSMAR重建基准名单和月收益，再用R展开公式复算17组实验指标。这比重复调用同一个函数多一层独立性；但两种语言共享原始资料和部分规则解释，仍可能共享数据偏差。本文据此支持的是既定月频模型具有可追溯定义且实现交叉一致，不声称已经消除原研报与本研究的一切差异。

% !TeX root = ../main.tex

\chapter{讨论与研究局限}
\section{四条路径的经济含义}
路径一保留持续成长与相对低估值的筛选骨架，以允许一次历史非正增长、相对排名和多维排序缓解候选稀缺。其全期收益高于E，但多项规则同步变化，不能把全部差异归于“19/20”或某一个因子。

路径二的价值在于恢复基础版可解释的全截面比较，同时展示排名分母如何进入组合结果；它在共同复现区间较优，在完整区间收益低于P1、回撤较小。规则更接近基础排名含义，不意味着历史业绩一定更高。

路径三、四保留各自选股名单，通过历史波动调整暴露，在当前数据中改善夏普和峰谷损失，但牺牲收益。P3全期指标较优，不应据此再叠加未验证的新规则、重新命名为唯一最优策略；四路径应共同保留，便于后续冻结检验。

\section{信息时点与参数选择}
最重要的限制是季度末财务映射。四路径的排名和风控修正均未解决财报实际披露时间和历史数值修订问题。此前冻结的公告敏感性结果显示，2003—2024年报告期映射年化25.40\%，完整公告约束为0.44\%；覆盖匹配仍为24.85\%。附录B保留其区间与定义，不将差值认定为前视偏差净因果效应，也不把该结果伪装成四路径已完成公告约束检验。

19/20及原有得分参数曾参考后续历史表现，P1与P2的保留选择也发生在研究者已见数据之后。追加区间是历史延长，不是严格事先冻结、从未接触的测试集。本文不报告未经校正的寻优显著性，也不把历史赢家排序表述为预测能力证据。

\section{收益、交易与统计边界}
收益采用含红利个股月收益与上证综指日收益复利，权重结构和分红口径并非完全一致。相对上证领先只是所设基准下的比较，不等于完成风险调整或证明可实现alpha。最大回撤基于月末，可能小于月内真实损失。

缺失收益再平均、月度重新等权、季度固定费用以及风控费用随仓位缩放，均是实现近似。未完整模拟涨跌停、停复牌、退市处置、冲击成本和融资约束。核心优先保留也不是当前组合逻辑。这些边界应在使用结果时与收益数字同时呈现。

\section{后续核验的优先顺序}
后续应先补证据，再增加策略复杂度。首先在CSMAR内部核查缺失公司的上市前比较期资料、合并报表类型和历史版本，分别记录“无报告期”“无同比基数”“基数非正”，避免将三者统称亏损。其次按原研报证券范围核对公告表和共同调仓日，只有事件时钟核对后，个股日行情验证才能针对真实差异。

若取得足够的CSMAR个股日行情，应保留原月频版本另建验证文件，在相同名单下检查调仓前后价格、分红复权、停复牌和持有收益；再比较月度等权与季度权重漂移。这样可以先判断收益实现差异，再判断选股差异，避免一次改变全部环节而无法解释原因。目前这些工作仍是待执行事项，正文没有将其写成已完成的证明。

最后才是在冻结规则下评价新增数据。全市场排名与候选池排名两条路径都应保留，风控分别与各自底层配对；任何核心优先、负PE排名移动或参数新设定都应另存文件、登记具体变化，不覆盖现有策略。这样才能继续保持本项目从基础到延展、再到优化的可读性和可审核性。

% !TeX root = ../main.tex

\chapter{结论}
\section{基础复现保留了什么，尚未复现什么}
本文以申万宏源柏顿·墨基尔成功选股法为对象，保留持续正成长与相对低估值两项筛选主线，并明确初始五年、成熟期历史、全市场排名、正PE和季度持有规则。基础复现B实际覆盖2003年4月至2015年6月147个月，年化38.86\%、累计收益5481.49\%、最大回撤$-$70.11\%。申万原文公布截至2015年7月31日年化36.74\%、累计收益4176.87\%、最大回撤幅度70.19\%。二者在收益数量级与深度回撤上相近，但并非相同日期、相同频率和相同名单的精确重建。

末期原文8只与B的9只仅重合4只，空仓时段也有差别。独立CSMAR诊断能够定位4只未匹配股票的历史同比和报告期缺口，原文发布日与固定季度末的调仓安排亦有明确差异。因此，合理的结论是完成了有明确近似边界的月频规则复现，而不是以年化接近宣布逐日一致。E年化26.17\%包含2015年后的长期经历，不能与原文36.74\%直接相减作为复现误差。

\section{四条路径的收益风险结论}
P1恢复正PE与初始五年基础规则，同时保留成熟期19/20、候选池排名、扩展候选与三维排序。全期年化29.64\%高于E，但最大回撤仍为$-$68.87\%，并未解决极端损失。其改善应归于当前整套规则，不能单独归功于某一个因子或将历史上更高收益等同于更合理的规则。

P2保持其他环节并恢复全截面两榜独立排名，年化27.38\%、最大回撤$-$64.15\%。它相对P1降低了全期年化与最深峰谷损失，却略微提高波动；原复现区间收益排序与全期相反。全市场排名的主要研究价值是恢复相对于基础版可解释的母体，同时揭示排名、补位及最终核心之间的机制，不能笼统称为全面优于P1。

P3、P4分别叠加相同仓位规则，年化降至23.19\%、20.85\%，最大回撤收窄至$-$32.74\%、$-$34.86\%。其风险改善伴随约6.5个百分点的全期年化代价，且零仓位月份均为0。已观察到的作用主要来自波动目标缩仓，不能声称25\%季度止损已经通过真实清仓得到验证；每一阶段的保护程度也不同。

原始时期补测还确认，147、148、151个月三种口径下，P1和P2的收益均高于基础版，P3和P4则低于基础版；151个月中两条选股路径分别提高年化2.33、4.52个百分点。单独考察2008年4月至2015年7月成熟期时，四路径收益均高于B，但受控路径仍低于各自未控底层。这些分组结果要求同时保留收益、风险、起止区间和配对对象，不能笼统声称四条路径在原始时期都提高绝对收益。

\section{持续性、程序证据与适用范围}
初始60个月共同的高复利明显影响全期结果，成熟期及后续阶段需要分开阅读。P1和P2的阶段排序变化，追加21个月中E又领先四路径，表明全期历史最优不是持续优势证明。完整年度计数、滚动五年和回撤恢复记录提供互补描述，均不是未经数据选择的独立样本外检验。

计算方面，B与E在147个共同月逐月一致，四路径初始20个季度256条名单与基础一致，全部正式持仓PE为正。独立基准重建E的560条股票—季度记录及282个月收益，最大月收益误差约$5\times10^{-16}$；R又对17组口径直接复算收益与回撤。上述证据支持既定模型的算术、连接和规则实现一致性，不能证明申万日频等价、原始数据绝对正确或策略具有可交易超额收益。

本研究因此共同保留B、E与P1—P4，而不根据一个终点淘汰其他路径。后续最优先事项是CSMAR历史资料及公告时钟核验、可成交收益和真实成本验证，再以冻结规则评价新增数据。本文的价值在于建立可读、可复核、能解释有利与不利结果的研究链条，并将收益追求、排名定义与风险暴露分别评价。

% !TeX root = ../main.tex

\chapter*{参考文献}
\addcontentsline{toc}{chapter}{参考文献}
\begin{thebibliography}{9}
\bibitem{ref1} 申万宏源研究. 申万大师系列·价值投资篇之十二：柏顿·墨基尔成功选股法[R]. 2015-10-14.
\bibitem{ref2} FAMA E F, FRENCH K R. The cross-section of expected stock returns[J]. The Journal of Finance, 1992, 47(2): 427-465. DOI: 10.1111/j.1540-6261.1992.tb04398.x.
\bibitem{ref3} ASNESS C S, MOSKOWITZ T J, PEDERSEN L H. Value and momentum everywhere[J]. The Journal of Finance, 2013, 68(3): 929-985. DOI: 10.1111/jofi.12021.
\end{thebibliography}

% !TeX root = ../main.tex

\appendix
\ctexset{chapter={name={附录},number={\Alph{chapter}},format+={\zihao{-2}\usefonttimes},afterskip={0.5\baselineskip},beforeskip={0.8\baselineskip}}}
\setcounter{chapter}{0}
\setcounter{table}{0}
\renewcommand{\thetable}{A\arabic{table}}
\chapter{年度收益完整表}
\begin{table}[htbp]
\centering
\caption{年度复利收益（\%）}
\renewcommand{\arraystretch}{1.18}
\footnotesize
\setlength{\tabcolsep}{3pt}
\begin{tabularx}{\textwidth}{@{}p{1.05cm}*{7}{C}@{}}
\toprule
年份 & 月数 & E & P1 & P2 & P3 & P4 & 上证 \\
\midrule
2003 & 9 & 36.32 & 36.32 & 36.32 & 26.16 & 26.16 & -0.90 \\
2004 & 12 & 6.84 & 6.84 & 6.84 & 4.79 & 4.79 & -15.40 \\
2005 & 12 & -0.05 & -0.05 & -0.05 & 0.33 & 0.33 & -8.33 \\
2006 & 12 & 126.39 & 126.39 & 126.39 & 84.78 & 84.78 & 130.43 \\
2007 & 12 & 446.28 & 446.28 & 446.28 & 182.09 & 182.09 & 96.66 \\
2008 & 12 & -65.02 & -60.67 & -52.44 & -25.19 & -20.23 & -65.39 \\
2009 & 12 & 91.75 & 169.65 & 143.31 & 61.08 & 60.70 & 79.98 \\
2010 & 12 & 39.48 & 31.00 & 39.52 & 18.98 & 27.37 & -12.15 \\
2011 & 12 & -9.15 & -17.17 & -20.47 & -18.11 & -20.82 & -21.68 \\
2012 & 12 & 20.58 & 21.73 & -2.02 & 19.08 & -3.07 & 3.17 \\
2013 & 12 & 37.33 & 19.28 & 37.24 & 13.93 & 23.70 & -6.75 \\
2014 & 12 & 49.37 & 50.14 & 53.96 & 51.88 & 46.43 & 52.87 \\
2015 & 12 & 38.87 & 62.56 & 71.35 & 54.29 & 71.35 & 8.49 \\
2016 & 12 & 30.82 & 5.14 & -3.22 & 1.36 & -1.46 & -12.31 \\
2017 & 12 & 5.18 & 14.10 & 6.97 & 13.01 & 6.69 & 6.56 \\
2018 & 12 & -6.41 & -25.06 & -27.19 & -25.06 & -27.19 & -24.59 \\
2019 & 12 & 55.91 & 34.89 & 34.50 & 33.01 & 32.15 & 22.30 \\
2020 & 12 & 17.81 & 25.10 & 24.40 & 25.19 & 24.88 & 13.87 \\
2021 & 12 & -17.86 & 48.86 & 58.92 & 48.82 & 55.89 & 4.80 \\
2022 & 12 & 13.95 & 8.08 & 21.89 & 3.94 & 13.22 & -15.13 \\
2023 & 12 & -8.77 & 9.89 & 4.54 & 9.09 & 1.12 & -3.70 \\
2024 & 12 & 29.52 & 38.52 & 17.09 & 38.12 & 16.90 & 12.67 \\
2025 & 12 & 48.35 & 48.82 & 30.44 & 44.26 & 30.33 & 18.41 \\
2026 & 9 & 15.50 & 9.99 & -7.71 & 8.35 & -7.20 & -3.19 \\
\bottomrule
\end{tabularx}
\par\smallskip
\begin{minipage}{\textwidth}\small 注：B与E在共同区间完全一致，不重复列B的自然年度短表。2003年为4至12月，2026年为1至9月；其余为12个月。均为复利实际收益，不将不足一年区间在此表内年化。\end{minipage}
\end{table}

\setcounter{table}{0}
\renewcommand{\thetable}{B\arabic{table}}
\chapter{历史冻结的信息时点诊断}
本附录仅使用此前已核验的261个月结果，不混入新追加21个月。旧基准月收益与本次E的前261个月逐月一致。固定滞后一季不是完整公告约束；完整约束需要同比分母、TTM依赖和历史成长记录等财务依赖均已公告，并使用当时可得的最新财务期。该计算仍使用最终会计值，不能称为完整还原历史版本的数据库。

\begin{table}[htbp]
\centering
\caption{历史冻结公告敏感性（仅261个月）}
\renewcommand{\arraystretch}{1.18}
\small
\setlength{\tabcolsep}{4.5pt}
\begin{tabularx}{\textwidth}{@{}p{3.15cm}*{3}{C}@{}}
\toprule
实施口径 & 年化收益/\% & 波动/\% & 最大回撤/\% \\
\midrule
报告期映射 & 25.40 & 32.50 & -70.11 \\
固定滞后一季 & 10.38 & 32.22 & -72.30 \\
完整公告约束 & 0.44 & 33.37 & -89.91 \\
仅当期公告 & 6.31 & 32.66 & -72.30 \\
公告覆盖匹配 & 24.85 & 31.78 & -68.06 \\
\bottomrule
\end{tabularx}
\end{table}

\begin{figure}[htbp]
\centering
\includegraphics[width=\linewidth]{fig16}
\caption{历史冻结区间的公告敏感性}
\par\smallskip
\begin{minipage}{0.97\linewidth}\small 注：仅2003-04至2024-12；不是四条优化路径的公告约束收益，不能外推为2026年统计。资料来源：CSMAR，作者计算。\end{minipage}
\end{figure}

覆盖匹配仍允许提前使用财务值，仅用于观察覆盖剔除能否解释差异。完整约束同时改变财务期和排名集合，不能用两者CAGR之差识别单一偏差的因果效应。此前HAC、因子回归和区块自助研究连同原报告保留在归档目录，本轮主结论不把其旧区间统计移植到四路径。

\setcounter{table}{0}
\renewcommand{\thetable}{C\arabic{table}}
\chapter{核验与复现记录}
基础与延展147个共同月收益一致；四路径初始20个季度256条名单与基础一致；P1与P2的历史候选池、三维得分一致；P2的217,083条成长率名次和263,221条PE名次与基础延展全榜一致；全部正式持仓PE为正。Python独立重算全部分段指标，并逐月重建P3、P4仓位及实际收益，与R结果在浮点精度内一致。

研究报告不附代码或原始数据。完整代码、输入版本校验、运行结果、十六组实证图和R环境信息见配套RMarkdown及其HTML；六方案结果和四条正式路径程序独立保存。原有策略程序和原始数据保持不变，修订前两份报告另作备份。

本机R保存结果并输出HTML后仍以0xC0000005退出，原因未定位。断言通过不等同于进程正常；日志及状态随附，Python独立复算和构建正常完成。

\setcounter{table}{0}
\renewcommand{\thetable}{D\arabic{table}}
\chapter{回撤恢复与完整年度的补充证据}
\section{最深月末回撤的发生与恢复}
下表从已有月收益复利计算，不生成新策略。峰值含收益开始前的净值1；同值峰顶取谷值之前最后一个月，恢复指谷值后首次达到该峰值。所列月数为两个观察月的间隔，不是日历交易日；只描述最深回撤这一段，不是所有回撤的平均时长，也不表示峰值至恢复期间始终持有同一名单。

\begin{table}[htbp]
\centering
\caption{最大月末回撤对应的时间与恢复}
\renewcommand{\arraystretch}{1.18}
\footnotesize
\setlength{\tabcolsep}{3pt}
\begin{tabularx}{\textwidth}{@{}p{1.15cm}*{6}{C}@{}}
\toprule
方案 & 峰值月份 & 谷值月份 & 回撤/\% & 至谷/月 & 首次恢复 & 至恢复/月 \\
\midrule
B & 2007-09 & 2008-10 & -70.11 & 13 & 2011-02 & 41 \\
E & 2007-09 & 2008-10 & -70.11 & 13 & 2011-02 & 41 \\
P1 & 2007-09 & 2008-10 & -68.87 & 13 & 2009-12 & 27 \\
P2 & 2007-09 & 2008-10 & -64.15 & 13 & 2009-11 & 26 \\
P3 & 2007-09 & 2008-10 & -32.74 & 13 & 2009-10 & 25 \\
P4 & 2015-05 & 2018-10 & -34.86 & 41 & 2020-06 & 61 \\
\bottomrule
\end{tabularx}
\end{table}

基准E最深回撤发生于2007年9月至2008年10月，2011年2月首次恢复此前峰值，峰值到恢复间隔41个月。P3最深回撤对应同一峰谷阶段，2009年10月恢复，间隔25个月；P4最深回撤则出现在2015年5月至2018年10月，到2020年6月恢复，共61个月。这说明较小最大回撤幅度不保证较短的恢复时间，且不同路径的最深回撤未必发生在同一事件中。

B的终点为2015年6月，其他方案为2026年9月，不能只比较“是否恢复”而忽略观察机会不同。月末统计不显示月内跌幅；原研报的日频峰谷日期不能用本表月份直接替代。每条最大回撤已与主表指标逐项复核。

\section{完整自然年度的收益分布}
\begin{table}[htbp]
\centering
\caption{2004—2025年完整年度对照}
\renewcommand{\arraystretch}{1.18}
\footnotesize
\setlength{\tabcolsep}{3pt}
\begin{tabularx}{\textwidth}{@{}p{1.15cm}*{5}{C}@{}}
\toprule
方案 & 年数 & 正收益年 & 领先E年 & 最差年 / \% & 最好年 / \% \\
\midrule
E & 22 & 16 & — & 2008 / -65.02 & 2007 / 446.28 \\
P1 & 22 & 18 & 11 & 2008 / -60.67 & 2007 / 446.28 \\
P2 & 22 & 16 & 10 & 2008 / -52.44 & 2007 / 446.28 \\
P3 & 22 & 19 & 9 & 2008 / -25.19 & 2007 / 182.09 \\
P4 & 22 & 17 & 7 & 2018 / -27.19 & 2007 / 182.09 \\
\bottomrule
\end{tabularx}
\end{table}

P1有18个正收益年度，高于E的16个，但22年中只有11年严格领先E；P3有19个正收益年度，领先E的年度仅9个。正收益年度更多与相对基准领先更多并非同一概念，绝对收益、相对收益和风险应分别阅读。P4最差完整年度为2018年，而其他路径均为2008年，和回撤时段不同的观察相呼应。

只纳入各12个月的22个年度；领先E指当年复利严格大于E，相等不计领先。E作为自身基准不报告领先次数。最好和最差仅是已观察年度极值，并非风险预测区间；跨年度次数没有控制共同市场暴露和参数选择，不当作显著性检验或未来胜率。

\end{document}
